% from ega0-1-fr.tex
\oldpage[0\textsubscript{I}]{70}

Pour tout voisinage $V$ de $0$ dans $A$, soit $V'$ l'ensemble des
$x=\sum_\alpha c_\alpha T^\alpha\in A'$ telles que $c_\alpha\in V$ pour tout
$\alpha$. On vérifie aussitôt que les $V'$ forment un système fondamental de
voisinages de $0$ définissant sur $A'$ une topologie d'anneau \emph{séparée}~;
nous allons définir canoniquement un \emph{isomorphisme topologique} de
l'anneau $A\{T_1,\ldots,T_r\}$ sur $A'$. Pour tout
$\alpha\in\mathbf{N}^r$ et tout $\lambda$, soit
$\varphi_{\lambda,\alpha}$ l'application de
$(A/\mathfrak J_\lambda)[T_1,\ldots,T_r]$ dans
$A/\mathfrak J_\lambda$ qui, à tout polynôme du premier anneau, fait
correspondre le coefficient de $T^\alpha$ dans ce polynôme. Il est clair que
les $\varphi_{\lambda,\alpha}$ forment un système projectif d'homomorphismes
de $(A/\mathfrak J_\lambda)$-modules, dont la limite projective est un
homomorphisme continu
$\varphi_\alpha:A\{T_1,\ldots,T_r\}\to A$~; nous allons voir que, pour tout
$y\in A\{T_1,\ldots,T_r\}$, la série formelle
$\sum_\alpha\varphi_\alpha(y)T^\alpha$ est \emph{restreinte}. En effet, si
$y_\lambda$ est la composante de $y$ dans $B_\lambda$, et si on désigne par
$H_\lambda$ l'ensemble fini des $\alpha\in\mathbf{N}^r$ pour lesquels les
coefficients du polynôme $y_\lambda$ ne sont pas nuls, on a
$\varphi_{\lambda,\alpha}(y_\mu)\in\mathfrak J_\lambda$ pour
$\mathfrak J_\mu\subset\mathfrak J_\lambda$ et $\alpha\notin H_\lambda$, et
par passage à la limite, $\varphi_\alpha(y)\in\mathfrak J_\lambda$ pour
$\alpha\notin H_\lambda$. On définit donc un homomorphisme d'anneaux
$\varphi:A\{T_1,\ldots,T_r\}\to A'$ en posant
$\varphi(y)=\sum_\alpha\varphi_\alpha(y)T^\alpha$, et il est immédiat que
$\varphi$ est continu. Inversement, si $\theta_\lambda$ est l'homomorphisme
canonique $A\to A/\mathfrak J_\lambda$, pour tout élément
$z=\sum_\alpha c_\alpha T^\alpha\in A'$ et tout $\lambda$, il n'y a qu'un
nombre fini d'indices $\alpha$ tels que $\theta_\lambda(c_\alpha)\neq0$, et
par suite
$\psi_\lambda(z)=\sum_\alpha\theta_\lambda(c_\alpha)T^\alpha$ appartient à
$B_\lambda$~; les $\psi_\lambda$ sont continus et forment un système
projectif d'homomorphismes dont la limite projective est un homomorphisme
continu $\psi:A'\to A\{T_1,\ldots,T_r\}$~; il reste enfin à vérifier que
$\varphi\circ\psi$ et $\psi\circ\varphi$ sont les automorphismes identiques,
ce qui est immédiat.
\end{env}

\begin{env}[7.5.2]
\label{0.7.5.2-fr}
Nous identifierons $A\{T_1,\ldots,T_r\}$ à l'anneau $A'$ des séries formelles
restreintes au moyen des isomorphismes définis dans (7.5.1). Les isomorphismes
canoniques
\[
  ((A/\mathfrak J_\lambda)[T_1,\ldots,T_r])
  [T_{r+1},\ldots,T_s]
  \xrightarrow{\sim}
  (A/\mathfrak J_\lambda)[T_1,\ldots,T_s]
\]
définissent, par passage à la limite projective, un isomorphisme canonique
\[
  \Bigl[(A\{T_1,\ldots,T_r\})\{T_{r+1},\ldots,T_s\}
  \xrightarrow{\sim}
  A\{T_1,\ldots,T_s\}
\]
\end{env}

\begin{env}[7.5.3]
\label{0.7.5.3-fr}
Pour tout homomorphisme continu $u:A\to B$ de $A$ dans un anneau linéairement
topologisé $B$, séparé et complet, et tout système $(b_1,\ldots,b_r)$ de $r$
éléments de $B$, il existe \emph{un homomorphisme continu, et un seul},
$\overline u:A\{T_1,\ldots,T_r\}\to B$, tel que $\overline u(a)=u(a)$ pour
tout $a\in A$ et $\overline u(T_j)=b_j$ pour $1\leq j\leq r$. Il suffit en
effet de prendre
\[
  \overline u\Bigl(\sum_\alpha c_\alpha T^\alpha\Bigr)
  =\sum_\alpha u(c_\alpha)b_1^{\alpha_1}\cdots b_r^{\alpha_r}~;
\]
les vérifications du fait que la famille
$(u(c_\alpha)b_1^{\alpha_1}\cdots b_r^{\alpha_r})$ est sommable dans $B$ et
de la continuité de $\overline u$ sont immédiates et laissées au lecteur. On
notera que cette propriété (pour $B$ et les $b_j$ arbitraires)
\emph{caractérise} l'anneau topologique $A\{T_1,\ldots,T_r\}$ à un
isomorphisme unique près.
\end{env}

\begin{proposition}[7.5.4]
\label{0.7.5.4-fr}
\begin{enumerate}
  \item[(i)] Si $A$ est un anneau admissible, il en est de même de
  $A'=A\{T_1,\ldots,T_r\}$.
  \item[(ii)] Soient $A$ un anneau adique, $\mathfrak J$ un idéal de définition
  de $A$ tel que $\mathfrak J/\mathfrak J^2$ soit de type fini

